\section*{Hoofdstuk \ref{ch:b2:reproductive-hormones} --- Hormonale regeling van de voortplanting}

\begin{solution}{exo:b2:reproductive-hormones:1}
Hypothalamus (GnRH, pulsen) $\to$ hypofysevoorkwab (LH, FSH) $\to$
\omterm{def:b2:mammal-reproduction:testis}{testis}: LH op de leydigcellen (testosteron), FSH op de sertolicellen
(\omterm{def:b2:mammal-reproduction:testis}{spermatogenese}, inhibine). Terugkoppeling: testosteron remt GnRH en
LH; inhibine remt FSH.
\end{solution}

\begin{solution}{exo:b2:reproductive-hormones:2}
Ovarium: \omterm{def:b2:reproductive-hormones:cycle}{folliculaire fase}, dag 1--14, oestradiol uit de groeiende
\omterm{def:b2:mammal-reproduction:ovary}{follikel}; \omterm{def:b2:mammal-reproduction:ovary}{ovulatie}, dag 14, \omterm{def:b2:reproductive-hormones:cycle}{LH-piek}; \omterm{def:b2:reproductive-hormones:cycle}{luteale fase}, dag 15--28,
progesteron uit het \omterm{def:b2:mammal-reproduction:ovary}{corpus luteum}. Baarmoeder: menstruatie, dag 1--5
(wegvallen van de steroïden); proliferatieve fase, dag 6--14 (oestradiol);
secretorische fase, dag 15--28 (progesteron).
\end{solution}

\begin{solution}{exo:b2:reproductive-hormones:3}
Het dagelijkse oestrogeen en progestageen houden LH en FSH door negatieve
terugkoppeling laag, zodat er geen \omterm{def:b2:mammal-reproduction:ovary}{follikel} wordt gerekruteerd, er geen
stijging van oestradiol optreedt en de piek nooit vuurt; het progestageen
maakt ook het slijm dikker. In de stopweek dalen de steroïden en wordt het
\omterm{def:b2:reproductive-hormones:cycle}{endometrium}, dat onder hun invloed is opgebouwd, afgestoten --- een
onttrekkingsbloeding, geen menstruatie na een \omterm{def:b2:mammal-reproduction:ovary}{ovulatie}.
\end{solution}

\begin{solution}{exo:b2:reproductive-hormones:4}
LH en FSH stijgen tot een veelvoud (geen terugkoppeling door testosteron
of inhibine); spieren, libido en de accessoire klieren gaan in regressie.
Testosteron herstelt de kenmerken en brengt LH weer omlaag, maar niet de
vruchtbaarheid: de testes zijn weg, en zelfs bij een man met testes
onderdrukt toegediend testosteron LH en daarmee het testosteron in de
\omterm{def:b2:mammal-reproduction:testis}{testis} dat de \omterm{def:b2:mammal-reproduction:testis}{spermatogenese} nodig heeft.
\end{solution}

\begin{solution}{exo:b2:reproductive-hormones:5}
$24/1.5 = 16$ pulsen per dag; $24/4 = 6$. Wanneer het \omterm{def:b2:mammal-reproduction:ovary}{corpus luteum} sterft,
zijn de pulsen nog traag, wat FSH bevoordeelt: FSH stijgt als eerste ---
het juiste hormoon, want FSH rekruteert het volgende cohort \omterm{def:b2:mammal-reproduction:ovary}{follikels}.
\end{solution}

\begin{solution}{exo:b2:reproductive-hormones:6}
\qty{200}{pg/mL} $= \qty{200}{ng/L}$: $200\times 10^{-9}/272 =
\qty{7.4e-10}{mol/L} = \qty{740}{pmol/L}$. \qty{15}{ng/mL} $=
\qty{15}{\micro g/L}$: $15\times 10^{-6}/314 = \qty{48}{nmol/L}$. In
één microliter: $7.4\times 10^{-16}$ mol $= 4.4\times 10^{8}$
moleculen.
\end{solution}

\begin{solution}{exo:b2:reproductive-hormones:7}
\omterm{def:b2:mammal-reproduction:ovary}{Ovulatie} op dag $35 - 14 = 21$; vruchtbaar van dag 18 tot dag 22. Bij
onregelmatige cycli kan de \omterm{def:b2:reproductive-hormones:cycle}{folliculaire fase}, en daarmee de ovulatiedag,
niet uit vorige cycli worden voorspeld, en wordt het venster gemist of
verkeerd geplaatst.
\end{solution}

\begin{solution}{exo:b2:reproductive-hormones:8}
Van dag 21 tot 26 zijn 2.5 verdubbelingen: $2^{2.5} = 5.7$ IU/L ---
genoeg. \omterm{def:b2:mammal-reproduction:implantation}{Innesteling} op dag 25 geeft $1.4$ IU/L op dag 26; het \omterm{def:b2:mammal-reproduction:ovary}{corpus
luteum}, dat al in regressie is, wordt niet gered, het progesteron daalt en
het embryo gaat met de menstruatie verloren. Late \omterm{def:b2:mammal-reproduction:implantation}{innesteling} is een
veelvoorkomende oorzaak van vroeg verlies.
\end{solution}

\begin{solution}{exo:b2:reproductive-hormones:9}
De hypofyse reageert op het patroon van GnRH, niet op de hoeveelheid:
pulsen om het uur houden LH en FSH in stand, dezelfde dagdosis continu
toegediend doet ze verdwijnen. Het sloot een eenvoudige dosis-responsrelatie
uit. Een maand langwerkende agonist: na een opflakkering van enkele dagen
daalt LH tot bijna nul (\omterm{prop:b2:reproductive-hormones:pulses}{desensitisatie van de receptoren}), oestradiol tot
postmenopauzale spiegels, het \omterm{def:b2:reproductive-hormones:cycle}{endometrium} atrofieert en er treedt geen
cyclus op --- een omkeerbare medicamenteuze menopauze.
\end{solution}

\begin{solution}{exo:b2:reproductive-hormones:10}
Geen hormoon: $R^* = 100/0.5 = 200$. Continu: $100/6.5 = 15$. Een
puls van twee minuten verwijdert $6\times 200\times(2/60) = 40$ receptoren
(\qty{20}{\%}); in het volgende uur krimpt het tekort met een factor
$e^{-0.5}$, d.w.z.\ \qty{39}{\%} ervan wordt hersteld. De voorraad stelt
zich in waar het verlies per puls gelijk is aan het herstel per uur:
ongeveer 150 receptoren vóór elke puls, drie kwart van het maximum ---
tien keer het continue niveau.
\end{solution}

\begin{solution}{exo:b2:reproductive-hormones:11}
Menopauze: geen \omterm{def:b2:mammal-reproduction:ovary}{follikels}, dus geen oestradiol of inhibine, dus geen
terugkoppeling: LH en FSH lopen hoog op. Prolactinoom: prolactine
vertraagt de GnRH-pulsgenerator, zodat LH laag is, de \omterm{def:b2:mammal-reproduction:ovary}{follikels} niet worden
aangestuurd, oestradiol laag is en er geen cyclus is. Granulosaceltumor:
constant hoog oestradiol onderdrukt LH door negatieve terugkoppeling (en
zonder rijpende \omterm{def:b2:mammal-reproduction:ovary}{follikel} valt er niets te ovuleren), zodat de cyclus
stopt bij een hoog oestradiol --- en het \omterm{def:b2:reproductive-hormones:cycle}{endometrium} overgroeit.
\end{solution}

\begin{solution}{exo:b2:reproductive-hormones:12}
Laag of kortstondig oestradiol verlaagt GnRH en LH; hoog, aanhoudend
oestradiol (boven \qty{200}{pg/mL} gedurende meer dan 36 uur) verandert de
respons van het kisspeptine-GnRH-systeem en van de hypofyse, die intussen
LH heeft opgeslagen, zodat hetzelfde hormoon nu een stoot uitlokt. De
\omterm{def:b2:mammal-reproduction:ovary}{ovulatie} moet alles-of-niets en getimed zijn: de eicel moet één keer
vrijkomen, op één dag, wanneer een \omterm{def:b2:mammal-reproduction:ovary}{follikel} rijp is. Een draaiknop (LH
evenredig met oestradiol) zou een geleidelijke, gedeeltelijke luteïnisatie
geven, \omterm{def:b2:mammal-reproduction:ovary}{follikels} die half rijp of met meerdere tegelijk ovuleren, en geen
welomschreven vruchtbare dag.
\end{solution}

\begin{solution}{pb:b2:reproductive-hormones:1}
\textbf{1.} Folliculair, dag 1--14: oestradiol stijgend, progesteron laag.
Luteaal, dag 15--28: progesteron hoog.
\textbf{2.} De piek bereikte haar maximum op dag 14 (begonnen op dag 13);
\omterm{def:b2:mammal-reproduction:ovary}{ovulatie} ongeveer 36 uur na het begin, op dag 15.
\textbf{3.} Dag 15 tot 28: 14 dagen, de levensduur van het \omterm{def:b2:mammal-reproduction:ovary}{corpus
luteum}.
\textbf{4.} Progesteron verhoogt de instelwaarde van de temperatuurregeling
met \qty{0.3}{\celsius} tot \qty{0.4}{\celsius}; de stijging verschijnt een
of twee dagen na de \omterm{def:b2:mammal-reproduction:ovary}{ovulatie}, zodat ze bevestigt dat de \omterm{def:b2:mammal-reproduction:ovary}{ovulatie} heeft
plaatsgevonden, maar haar niet kan aankondigen.
\textbf{5.} Dag 11 tot 14 (210, 260, 310, 220): drie tot vier dagen, meer
dan de 36 uur die de omslag vereist.
\textbf{6.} Dag 12 tot 16.
\textbf{7.} $310\times 10^{-9}/272 = \qty{1.14e-9}{mol/L} =
\qty{1140}{pmol/L}$; $16\times 10^{-6}/314 = \qty{51}{nmol/L}$.
\textbf{8.} $65/8 = 8$. Met een halveringstijd van een uur halveert LH elk
uur zodra de secretie stopt, zodat het een dag later weer bijna op het
basisniveau is.
\textbf{9.} \omterm{def:b2:mammal-reproduction:ovary}{Ovulatie} rond dag 22, een cyclus van ongeveer 36 dagen.
\textbf{10.} \omterm{def:b2:mammal-reproduction:implantation}{Innesteling} op dag 22; tegen dag 27 (vijf dagen, 2.5
verdubbelingen) ongeveer \qty{5.7}{IU/L}.
\textbf{11.} Een progesteron dat op dag 27 tot \qty{2}{ng/mL} is gedaald,
betekent dat het \omterm{def:b2:mammal-reproduction:ovary}{corpus luteum} volgens schema in regressie gaat: niets heeft
het gered.
\textbf{12.} Progesteron rond \qty{20}{ng/mL} en stijgend, oestradiol rond
\qty{200}{pg/mL}.
\textbf{13.} Aan het eind van de cyclus dalen oestradiol, progesteron en
inhibine samen en ontsnapt FSH aan zijn terugkoppeling; naarmate het cohort
groeit, stijgen oestradiol en inhibine en daalt FSH, en alleen de \omterm{def:b2:mammal-reproduction:ovary}{follikel}
met de meeste FSH-receptoren (de meeste granulosacellen) blijft op het
dalende FSH groeien --- de rest sterft.
\textbf{14.} $R^* = 200$ zonder hormoon; $100/6.5 = 15$ onder
continu hormoon.
\textbf{15.} $6\times 200\times 2/60 = 40$ receptoren, \qty{20}{\%};
tijdconstante $1/d = \qty{2}{h}$, zodat in een uur \qty{39}{\%} van het
tekort wordt hersteld; de voorraad stelt zich in rond 150 receptoren, drie
kwart van het maximum.
\textbf{16.} Pulsen houden de hypofyse op ongeveer 150 receptoren, en ze
beantwoordt elke puls met LH; een continu infuus drijft de voorraad naar 15
en dezelfde cellen vallen stil, ongeacht de dosis; pulsen herstellen de
voorraad en de respons.
\textbf{17.} Eerste dag: LH en testosteron stijgen (de agonist stimuleert
receptoren die er nog zijn --- de opflakkering). Na een maand: receptoren
gedesensitiseerd, LH bijna nul, testosteron op castratieniveau, wat het
therapeutische doel is.
\textbf{18.} FSH stijgt (geen inhibine), LH blijft normaal; elke cyclus
overleven meer \omterm{def:b2:mammal-reproduction:ovary}{follikels} de selectie, met meervoudige \omterm{def:b2:mammal-reproduction:ovary}{ovulaties} en
twee-eiige tweelingen.
\textbf{19.} In een natuurlijke cyclus doodt de daling van FSH alle
\omterm{def:b2:mammal-reproduction:ovary}{follikels} behalve de dominante; een constante toevoer van FSH laat het hele
cohort doorgroeien tot rijpheid.
\textbf{20.} Tien \omterm{def:b2:mammal-reproduction:ovary}{follikels} maken al vroeg veel meer dan \qty{200}{pg/mL}
oestradiol, wat een voortijdige \omterm{def:b2:reproductive-hormones:cycle}{LH-piek} en \omterm{def:b2:mammal-reproduction:ovary}{ovulatie} vóór de oogst zou
uitlokken; de antagonist blokkeert GnRH, zodat er geen piek kan optreden.
\textbf{21.} In de tabel piekte LH op dag 14 en volgde de \omterm{def:b2:mammal-reproduction:ovary}{ovulatie} ongeveer
36 uur na het begin van de piek; hCG vervangt de piek, en een oogst na 35
uur vangt de oöcyten rijp maar nog in de \omterm{def:b2:mammal-reproduction:ovary}{follikel}.
\textbf{22.} $10\times 0.7 = 7$ embryo's, $7\times 0.4 = 2.8$
\omterm{def:b2:mammal-reproduction:implantation}{blastocysten}; $1 - 0.65^{2} = 0.58$.
\textbf{23.} Constant oestrogeen en progestageen houden FSH en LH laag: er
groeit geen \omterm{def:b2:mammal-reproduction:ovary}{follikel}, oestradiol stijgt nooit tot de drempel, geen piek,
geen \omterm{def:b2:mammal-reproduction:ovary}{corpus luteum}, geen stijging van progesteron.
\textbf{24.} Testosteron moet worden gegeven voor de secundaire kenmerken
en het libido; maar het herstelt niet de hoge concentratie in de \omterm{def:b2:mammal-reproduction:testis}{testis} die
de leydigcellen leverden, zodat de \omterm{def:b2:mammal-reproduction:testis}{spermatogenese} onderdrukt blijft --- wat
precies de bedoeling is.
\textbf{25.} \omterm{def:b2:mammal-reproduction:ovary}{Ovulatie} op dag 15; \omterm{def:b2:reproductive-hormones:cycle}{luteale fase} 14 dagen; vruchtbaar venster
dag 12--16; piek uitgelokt door oestradiol boven \qty{200}{pg/mL}
gedurende meer dan 36 uur, bereikt op dag 11--14.
\end{solution}
